\chapter{Come questo programma decide che qualcosa è vero}
\label{ch:method}

Un codice di progetto vale quanto le prove che è corretto. Questo
capitolo enuncia le due regole cui ogni suite di test del documento
obbedisce, spiega perché sono state scelte e definisce i pochi
termini usati per riportare i risultati. Si chiude con la tabella
dei simboli valida per l'intero documento.

\section{Regola uno: nessuna tolleranza è scelta a mano}

Ogni verifica di questo documento confronta un numero calcolato con
una banda, e \emph{nessuna banda è scelta da un essere umano}. La
ragione è semplice: una tolleranza scelta dopo aver visto la
risposta non certifica nulla. Esiste sempre, e passa sempre.

La banda è invece \emph{derivata}, e in quasi tutti i casi è il
calcolo stesso a derivarla su di sé. La stessa quantità viene
calcolata due volte, una su una discretizzazione grossolana e una su
una fine. La differenza fra le due è una misura dell'errore che lo
schema sta ancora commettendo. La banda è quella differenza,
moltiplicata per un fattore di sicurezza fisso
$K_{\rm RICH} = 4$, più un pavimento al livello
dell'arrotondamento in virgola mobile:
\begin{equation}
\text{banda} \;=\; K_{\rm RICH}\,
\bigl|\,\text{grossolana} - \text{fine}\,\bigr|
\;+\; \text{pavimento di arrotondamento}.
\label{eq:band}
\end{equation}
È la classica idea dell'\emph{estrapolazione di Richardson} usata
come stima dell'errore anziché come acceleratore, e il numero che ne
risulta è chiamato ovunque \emph{banda di Richardson}
(Figura~\ref{fig:method}b). Poiché il metodo è del secondo ordine
(Capitolo~\ref{ch:moc}), dimezzare il passo dovrebbe ridurre
l'errore di circa quattro volte; quando una suite riporta un
``rapporto grossolana/fine di 4.52'' sta riportando che lo schema
converge al ritmo che dovrebbe, il che è già di per sé una prova
che la banda significa qualcosa.

Non tutte le quantità hanno una discretizzazione da raffinare. Dove
non c'è, la banda è derivata da qualcos'altro di misurabile --- la
dispersione dei dati di riferimento stessi, o la scala fisica della
quantità --- e la fonte è sempre dichiarata. Un capitolo (il
Capitolo~\ref{ch:genotwin}) usa soglie \emph{dichiarate} anziché
derivate; lo dice esplicitamente, ne dà la provenienza e spiega
perché il raffinamento sarebbe lì il metro sbagliato. Quell'eccezione
è deliberata e visibile, che è appunto il punto.

\section{Regola due: ogni suite deve saper dire di no}

Un test che non ha mai fallito potrebbe non misurare nulla. Perciò,
accanto alle verifiche positive, ogni suite contiene corruzioni
deliberate --- chiamate \emph{rigettatori} --- che la suite è tenuta
a cogliere: un segno invertito, un termine eliminato, una pressione
ambiente sbagliata, un piede caratteristico sbagliato, una geometria
perturbata di una frazione di percento. Se un rigettatore \emph{non}
fallisce, la verifica positiva corrispondente non vale nulla e il
passo non conta come fatto.

I rigettatori fanno lavoro reale in questo documento. Sono ciò che
rende affidabile una soglia dichiarata (mostrano che lo strumento
risolve difetti di un ordine di grandezza superiori al divario che
si sta misurando), e parecchi di essi hanno colto errori genuini nel
programma prima che potessero raggiungere un risultato.

La Figura~\ref{fig:method}a mette entrambe le regole in un'unica
immagine: ogni verifica del documento, espressa come deviazione
misurata divisa per la propria banda, insieme a ogni corruzione
deliberata. Le verifiche positive devono cadere a sinistra di uno; i
rigettatori devono cadere a destra. Nulla è vicino alla linea per
caso.

\begin{figure}[htbp]
\centering
\includegraphics[width=\textwidth]{fig_method.pdf}
\caption{Il metodo di questo documento, in un'immagine. (a)~Ogni
verifica certificata (blu) e ogni corruzione deliberata (rosso),
riportate come deviazione misurata divisa per la banda derivata
\emph{propria} di quella verifica --- così l'asse orizzontale è
adimensionale e la linea verticale a uno è la soglia di
superamento/fallimento per tutte in una volta. Le verifiche positive
cadono molto a sinistra (la maggior parte da due a sei ordini di
grandezza); i rigettatori cadono molto a destra. (b)~Da dove viene
una banda: la stessa quantità è calcolata su una griglia grossolana
e su una fine, la differenza misura l'errore che lo schema sta
ancora commettendo, e la banda è quella differenza per un fattore di
sicurezza quattro più un pavimento di arrotondamento. Gli assi sono
la spaziatura di griglia $h$ e il residuo della verifica, entrambi
relativi e logaritmici; la linea tratteggiata è la pendenza del
secondo ordine che il metodo dovrebbe seguire.}
\label{fig:method}
\end{figure}

\section{Come vengono riportati i risultati}

Ricorrono alcune convenzioni.

\emph{``$m/n$ PASS''} conta le verifiche di una suite, rigettatori
compresi: una suite con sei verifiche positive e due corruzioni che
si comporta correttamente riporta 8/8.

\emph{``Certificato da Newton''} significa che il residuo convergente
di ogni cella è stato verificato essere al livello
dell'arrotondamento, scalato con la dimensione delle sue incognite
(Capitolo~\ref{ch:moc}). Il numero citato (``residuo peggiore
0.019'') è una frazione della tolleranza ammessa, quindi qualunque
valore sotto uno passa.

\emph{``Risultato di riferimento''} significa che il numero è citato
alla lettera da un'esecuzione reale del carrier nominato, non
ricalcolato per il documento. Il Capitolo~\ref{ch:status} elenca
l'unico comando che riproduce ciascuno.

\emph{Un oracolo} è un caso la cui risposta è nota indipendentemente
dal programma --- una soluzione in forma chiusa, un'identità esatta
di conservazione o un secondo codice --- rispetto al quale il
programma può essere giudicato. Dove non esiste un oracolo, una
suite ripiega sulle \emph{chiusure interne}: affermazioni di
conservazione (massa, quantità di moto) che devono valere qualunque
sia la risposta, più i \emph{controlli di vitalità}, che verificano
che un termine che il programma dichiara di includere cambi
effettivamente la risposta quando viene rimosso.

Infine, e deliberatamente: \textbf{in questo programma anche i
difetti colti dai test sono risultati.} I capitoli principali
descrivono ciascun passo così come funziona oggi; ogni trappola,
ogni svolta sbagliata e ogni lezione strutturale misurata è raccolta
nell'Appendice~\ref{ch:appendix}, con riferimenti incrociati dal
passo cui appartiene. Un difetto colto è la prova che lo strumento
ha i denti.

\section{Simboli usati nel documento}

\begin{table}[htbp]
\centering
\small
\begin{tabular}{@{}llp{7.4cm}@{}}
\toprule
Simbolo & Unità & Significato \\
\midrule
\multicolumn{3}{@{}l}{\emph{Geometria e flusso}}\\
$x$, $y$ & m & posizione assiale, raggio \\
$u$, $v$ & m/s & componenti assiale e radiale della velocità \\
$q$ & m/s & modulo della velocità, $\sqrt{u^2+v^2}$ (indicato $W$
  negli articoli di Rao) \\
$\theta$ & rad, gradi & direzione del flusso, $\arctan(v/u)$ \\
$p$, $\rho$, $T$ & Pa, kg/m$^3$, K & pressione statica, densità,
  temperatura \\
$c$ & m/s & velocità locale del suono \\
$M$ & -- & numero di Mach, $q/c$ \\
$\alpha$ & rad & angolo di Mach, $\sin\alpha = 1/M$ \\
$C^{+}$, $C^{-}$ & -- & le due famiglie caratteristiche, a
  $\theta+\alpha$ e $\theta-\alpha$ \\
$\nu(M)$ & rad & funzione di Prandtl--Meyer (la rotazione prodotta
  da un'espansione di spigolo) \\
$\delta$ & -- & 0 per flusso piano, 1 per assialsimmetrico \\
\midrule
\multicolumn{3}{@{}l}{\emph{La macchina e il suo punto di
  funzionamento}}\\
$y_t$ & m & raggio di gola (tenerlo fisso fissa la portata) \\
$r_{tu}$, $r_{td}$ & m & raggi di raccordo in gola, a monte e a
  valle \\
$A_t$, $A_e$ & m$^2$ & area di gola e di uscita \\
$\varepsilon$ & -- & rapporto di espansione, $A_e/A_t$ \\
$L$ & m & lunghezza della parete \\
$P_0$, $T_0$ & Pa, K & pressione e temperatura di camera (di
  ristagno) \\
$p_a$ & Pa & pressione ambiente (di contropressione) \\
$p_e$, $u_e$ & Pa, m/s & pressione e velocità di uscita \\
$\dot m$ & kg/s & portata \\
$J$, $F$ & N & spinta (l'obiettivo da ottimizzare) \\
\midrule
\multicolumn{3}{@{}l}{\emph{Il ciclo (motore a detonazione
  rotante)}}\\
$\xi$ & -- & fase del ciclo, $0$ al passaggio dell'onda e $1$ alla
  fine dello svuotamento \\
$\mu$ & -- & la misura che pesa le fasi nella media di ciclo \\
$PR$, $TR$ & -- & rapporto di pressione e di temperatura sul
  ciclo \\
\midrule
\multicolumn{3}{@{}l}{\emph{Ugello a spina e swirl}}\\
$l$ & m & lunghezza di troncamento: dove la spina viene tagliata \\
$D$ & -- & il punto di troncamento stesso \\
$w$ & m/s & velocità azimutale (di swirl) \\
$\Gamma$ & m$^2$/s & circolazione, $y\,w$; costante in un vortice
  libero \\
\midrule
\multicolumn{3}{@{}l}{\emph{Strumenti}}\\
$R_{\rm mom}$ & N & disaccordo fra le due vie indipendenti di
  calcolo della spinta \\
$R_{\rm corner}$ & Pa & residuo terminale (di spigolo): nullo a un
  labbro ottimale \\
$f_2$ & m/s & primo integrale di Rao,
  $q\cos(\theta-\alpha)/\cos\alpha$, costante lungo la
  caratteristica terminale di un ugello ottimale \\
$\lambda_2$, $\lambda_3$ & -- & moltiplicatori di Lagrange costanti
  di Rao ($f_2 = -\lambda_2$) \\
$K_{\rm RICH}$ & -- & fattore di sicurezza di Richardson, fissato
  a 4 \\
\bottomrule
\end{tabular}
\caption{Simboli usati in tutto questo documento. Le lunghezze sono
espresse in metri; i casi di riferimento sono posti con raggio di
gola $y_t = 1$~m, quindi i valori numerici di $x$ e $y$ sono anche
multipli del raggio di gola.}
\label{tab:notation}
\end{table}
